\chapter{Distribuciones de campo hiperfino}
\label{ch:distribuciones}

Si ha llegado hasta aquí sabe reconocer y ajustar un singlete, un doblete o un
sexteto: sabe que cada uno corresponde a un \textbf{entorno de hierro bien
definido}, con un único desplazamiento isomérico, un único cuadrupolo y —en el
sexteto— un único campo hiperfino. Las \emph{distribuciones} son el paso
siguiente, y su lógica es distinta. Este capítulo empieza \textbf{por la física}:
qué son, de dónde salen y qué significa cada parámetro. Solo después entra en los
controles del programa.

\section{La física: por qué aparecen las distribuciones}
\label{sec:fisica-dist}

\subsection{Del cristal perfecto al material real}

En un \textbf{cristal perfecto} todos los átomos de Fe son equivalentes: ven el
mismo entorno, así que tienen \emph{exactamente} el mismo campo hiperfino
$\BHF$, el mismo $\delta$ y el mismo $\dEQ$. El resultado es un sexteto (o
doblete, o singlete) \textbf{discreto} y nítido, de líneas finas. Es el mundo de
los capítulos anteriores.

Pero la mayoría de los materiales \textbf{no son cristales perfectos}. En un vidrio
metálico, una aleación desordenada, un óxido no estequiométrico o una
nanopartícula, cada átomo de Fe ve un entorno \textbf{ligeramente distinto} del de
su vecino: puede tener otro número de átomos de un cierto tipo alrededor, otra
distancia de enlace, otra deformación local. Y como los parámetros hiperfinos
\emph{dependen} de ese entorno, dejan de tener un único valor: se
\textbf{reparten} en un rango.

\begin{nota}
La idea clave: en un material desordenado no hay «un» campo hiperfino, sino
\textbf{tantos campos como entornos locales distintos hay}. El espectro deja de
ser un sexteto y pasa a ser la \textbf{superposición de muchísimos sextetes},
cada uno con su campo, pesados por cuántos átomos de Fe tienen ese campo.
\end{nota}

\subsection{Qué es $P(\BHF)$}

Esa «cuántos átomos tienen cada campo» es precisamente la distribución. $P(\BHF)$
es una \textbf{densidad de probabilidad}: $P(\BHF)\,d\BHF$ es la fracción de
átomos de Fe cuyo campo hiperfino cae entre $\BHF$ y $\BHF + d\BHF$. Dicho de otro
modo, es el \textbf{histograma de los entornos locales} vistos por el hierro.

\begin{itemize}[nosep]
  \item Un material \textbf{cristalino} tiene $P(\BHF)$ \emph{muy estrecha} (casi
        una delta): casi todos los átomos comparten el mismo campo.
  \item Un material \textbf{desordenado} tiene $P(\BHF)$ \emph{ancha}: los campos
        se reparten en un intervalo amplio.
  \item Un material con \textbf{dos poblaciones} (p.\,ej.\ dos fases, o núcleo y
        superficie de una nanopartícula) tiene $P(\BHF)$ con \textbf{dos máximos}.
\end{itemize}

El espectro medido es entonces
\[
  T(v) \;=\; b \;-\; \int P(\BHF)\; S(v;\BHF)\, d\BHF,
\]
la suma (integral) de todos los subespectros $S(v;\BHF)$ pesados por $P(\BHF)$.
Ajustar una distribución es \textbf{el problema inverso}: partir del espectro y
recuperar la $P(\BHF)$ que lo generó.

\begin{nota}
Por eso una distribución \emph{no se simula} de la nada: se \textbf{infiere del
espectro} (ajuste). Lo que sí puede simular es el \emph{espectro} que produciría
una $P(\BHF)$ que usted \emph{suponga} —útil para desarrollar intuición
(§\ref{sec:sim-cristalina})—. El desarrollo matemático del problema inverso está
en el anexo~\ref{ap:mates}.
\end{nota}

\subsection{Distribución no es lo mismo que línea ancha}

Este punto confunde a menudo y conviene fijarlo. Hay \textbf{dos maneras} de que
las líneas de un espectro salgan anchas, y significan cosas físicas distintas:

\begin{itemize}[nosep]
  \item \textbf{Anchura de línea $\Gamma$}: es el ensanchamiento \emph{intrínseco}
        de \emph{un} entorno (vida media del estado, efectos instrumentales).
        Ensancha \textbf{todas} las líneas por igual y de forma simétrica, sin
        cambiar sus posiciones.
  \item \textbf{Anchura de la distribución}: es el \emph{abanico de entornos}
        distintos. Como la posición de las líneas del sexteto escala con el campo
        (§\ref{sec:modelo-mates}), una dispersión de campos ensancha
        \textbf{más las líneas externas que las internas}, y puede producir
        formas \textbf{asimétricas}.
\end{itemize}

\begin{aviso}
No intente «tapar» una distribución subiendo $\Gamma$. Un sexteto único con
$\Gamma$ grande y una verdadera $P(\BHF)$ ancha \textbf{no} dan el mismo espectro:
la distribución ensancha las líneas externas más que las internas y puede ser
asimétrica, cosa que $\Gamma$ no reproduce. Si al ajustar de forma discreta las
anchuras salen «infladas» y las líneas externas peor ajustadas que las internas,
es señal de que hay una distribución.
\end{aviso}

\subsection{Qué le dice cada rasgo de la distribución}

Una vez obtenida $P(\BHF)$, su \textbf{forma} es información física:

\begin{center}
\begin{tabular}{@{}lp{0.62\textwidth}@{}}
\toprule
\textbf{Rasgo de $P(\BHF)$} & \textbf{Qué implica físicamente} \\
\midrule
\textbf{Posición} (media / máximo) & El campo \emph{típico}: refleja el entorno
medio (composición, coordinación, estado magnético medio). \\[3pt]
\textbf{Anchura} & El \textbf{grado de desorden}: cuanto más ancha, más varían los
entornos (más dispersión química o estructural). \\[3pt]
\textbf{Asimetría} (cola) & Un gradiente de entornos: p.\,ej.\ una minoría de
átomos con muchos vecinos de otro tipo genera una cola hacia campos altos o bajos. \\[3pt]
\textbf{Varios máximos} & \textbf{Subpoblaciones} distintas: dos fases, dos sitios
cristalográficos, núcleo/superficie de nanopartícula, etc. \\[3pt]
\textbf{Peso en $\BHF\!\approx\!0$} & Fracción \textbf{no magnética / relajante}
(paramagnética o superparamagnética a esa temperatura). \\
\bottomrule
\end{tabular}
\end{center}

El mismo razonamiento vale para $P(\dEQ)$ (dispersión de gradientes de campo
eléctrico en materiales paramagnéticos) y $P(\mathrm{IS})$ (dispersión de
densidades electrónicas / estados de oxidación), tratadas en el
anexo~\ref{ap:dist-is}.

\section{Idea: de un sexteto a una distribución}

Una distribución reparte la probabilidad entre una malla de campos
$\{\BHF^{(j)}\}$ con pesos $P_j \ge 0$. El espectro es
\[
  T(v) \;=\; b \;-\; \sum_j P_j \, S_j(v),
\]
donde $S_j$ es el sexteto correspondiente al campo $\BHF^{(j)}$. Un
\textbf{material cristalino} con un único campo se representa como una
distribución \emph{estrecha} (casi una delta) en torno a ese valor; un material
desordenado, como una distribución \emph{ancha}.

\autofig[\textwidth]{fig_pbhf_distribution}{Izquierda: distribución de campo
hiperfino cristalina (nítida, en torno a 49~T) frente a una distribución ancha.
Derecha: los espectros simulados correspondientes. La distribución ancha
«emborrona» y solapa las líneas del sexteto.}

Observe en la figura cómo un $P(\BHF)$ estrecho produce un sexteto de líneas
finas (comportamiento cristalino), mientras que uno ancho da un espectro con
líneas ensanchadas y solapadas.

\section{Simular el espectro de un caso cristalino}
\label{sec:sim-cristalina}

El caso cristalino (un campo bien definido) es el punto de partida natural para
entender el modelo. En modo distribución:

\begin{enumerate}[nosep]
  \item Active el \textbf{modo distribución} y elija la variable ($P(\BHF)$ o
        $P(\dEQ)$).
  \item Fije el rango de la malla ($\BHF$ mínimo y máximo) y el número de bins.
  \item Fije los parámetros compartidos del kernel: $\delta$, $\dEQ$ y la anchura
        $\Gamma$ de línea.
  \item Suponga una distribución \emph{estrecha} centrada en el campo esperado
        (casi una delta): el espectro simulado tenderá al de un material
        cristalino, con líneas finas.
\end{enumerate}

El gráfico muestra a la vez la \textbf{distribución} $P(\BHF)$ supuesta y el
\textbf{espectro} que produce, de modo que se ve el efecto de estrechar o
ensanchar la distribución. Sobre datos reales, en cambio, se \textbf{ajusta}
$P(\BHF)$ (siguientes secciones).

\guicap{gui_distribution_panel}{Panel de distribución: variable, forma,
regularizador, rango de la malla y número de bins.}

\section{El panel de distribución, control a control}
\label{sec:panel-dist}

El panel de distribución reúne bastantes controles. Aquí se explica \textbf{cada
uno}: qué es, para qué sirve, qué valores son razonables y cuándo conviene
tocarlo. Muchos se \textbf{ocultan o deshabilitan} según la forma elegida (por
ejemplo, el regularizador solo aparece con Histograma), así que no se asuste si no
ve todos a la vez.

\begin{nota}
Idea de fondo: el ajuste de distribución construye \textbf{muchos subespectros}
(sextetes o dobletes), uno por cada valor de la variable en una malla, y busca
\textbf{qué peso} $P_j\ge 0$ tiene cada uno para reproducir el espectro. Los
parámetros del panel controlan (a) \emph{qué} se distribuye, (b) \emph{cómo} es
cada subespectro, (c) la \emph{malla} de valores y (d) la \emph{regularización}.
\end{nota}

\subsection{Qué se distribuye: la variable}
\label{sec:dist-variable}

Lo primero es decidir \textbf{qué magnitud} varía de un entorno a otro. Se elige
en el selector de modo:

\begin{center}
\begin{tabular}{@{}lp{0.66\textwidth}@{}}
\toprule
\textbf{Modo} & \textbf{Qué distribuye} \\
\midrule
$P(\BHF)$ & El \textbf{campo hiperfino} (materiales magnéticos: cada entorno
tiene su campo). Es el caso más común. \\[3pt]
$P(\dEQ)$ & El \textbf{desdoblamiento cuadrupolar} (materiales paramagnéticos:
distribución de gradientes de campo eléctrico, con dobletes en vez de sextetes). \\[3pt]
$P(\text{IS})$ & El \textbf{desplazamiento isomérico}. \\[3pt]
\textbf{2D} & Distribución conjunta de dos variables a la vez, p.\,ej.\
$P(\BHF,\dEQ)$ (§\ref{sec:dist-2d}). \\
\bottomrule
\end{tabular}
\end{center}

La variable distribuida es la que barre la malla; las \emph{demás} magnitudes se
fijan como parámetros globales del núcleo (siguiente apartado).

\subsection{Forma y regularizador}

\begin{itemize}[nosep]
  \item \menu{Forma}: la familia de $P$ que se ajusta —Histograma, Gaussiana,
        VBF, Binomial, Fija o 2D— (§\ref{sec:formas-dist} las detalla). Es la
        decisión que más condiciona el resto de controles.
  \item \menu{Regularizador}: solo con \textbf{Histograma}. Elige cómo se penaliza
        la falta de suavidad —\emph{tikhonov}, \emph{tv} o \emph{maxent}—
        (§\ref{sec:regularizacion}).
  \item \menu{Componentes VBF}: solo con \textbf{VBF}. Número de gaussianas (1--6)
        que suman la distribución. Empiece por 1--2 y añada solo si el ajuste lo
        pide.
\end{itemize}

\subsection{Parámetros del núcleo (cada subespectro)}
\label{sec:panel-kernel}

Cada subespectro de la malla comparte estos parámetros «de línea». Son los que
definen \emph{cómo} es un sexteto/doblete individual antes de repartir los pesos:

\begin{center}
\begin{tabular}{@{}lp{0.50\textwidth}ll@{}}
\toprule
\textbf{Control} & \textbf{Función} & \textbf{Típico} & \textbf{Rango} \\
\midrule
$\delta$ & Desplazamiento isomérico común a todos los subespectros (mueve el
patrón). Fíjelo al $\delta$ de su fase o déjelo libre. & $0$ & $-2{,}5\ldots2{,}5$ \\[4pt]
$\dEQ$ & Cuadrupolo global (en modo $P(\BHF)$). Se \textbf{deshabilita} si es la
variable distribuida. & $0$ & $-4\ldots4$ \\[4pt]
BHF fijo & Campo fijo usado cuando se distribuye $\dEQ$ o IS (no se usa en modo
$P(\BHF)$). & $33$ & $0\ldots60$ \\[4pt]
$\Gamma$ & \textbf{Anchura} de línea de cada subespectro. \emph{Clave}: controla
la resolución (véase el aviso). & $0{,}30$ & $0{,}06\ldots2{,}0$ \\
\bottomrule
\end{tabular}
\end{center}

\begin{aviso}
La anchura $\Gamma$ y el número de bins están \textbf{acoplados} con la
regularización. Si pone $\Gamma$ muy pequeña, cada subespectro es muy estrecho y
el histograma tiende a «picos» y a seguir el ruido; si la pone grande, la
distribución se ve más lisa pero pierde detalle. Un valor físico razonable
($\Gamma \approx 0{,}25$--$0{,}35$~mm/s para $^{57}$Fe) y una $\alpha$ elegida con
la L-curve dan el mejor compromiso.
\end{aviso}

\subsection{Correlación con el campo: $\kappa_\delta$ y $\kappa_q$}

Dos pendientes opcionales que acoplan $\delta$ y $\dEQ$ al valor del campo
(§\ref{sec:correlacion}). Por defecto valen \textbf{0} (sin correlación) y están
\textbf{fijas}; para refinarlas, desmarque su candado. Solo aplican a Histograma y
VBF.

\subsection{La malla de valores}

Define sobre qué valores de la variable se reparte la probabilidad:

\begin{center}
\begin{tabular}{@{}lp{0.58\textwidth}ll@{}}
\toprule
\textbf{Control} & \textbf{Función} & \textbf{Típico} & \textbf{Rango} \\
\midrule
mín & Extremo \textbf{inferior} de la malla. & $0$ & $0\ldots60$ \\[3pt]
máx & Extremo \textbf{superior}. Ajústelo al rango donde \emph{espera} campos
(p.\,ej.\ $0$--$55$~T). & $50$ & $0\ldots60$ \\[3pt]
bins & \textbf{Número de puntos} de la malla. Más bins = más resolución pero
problema más mal condicionado (más regularización necesaria). & $50$ & $10\ldots200$ \\
\bottomrule
\end{tabular}
\end{center}

\begin{truco}
Elija el \textbf{rango} lo más ceñido posible a donde hay señal: una malla que se
extiende por zonas vacías «gasta» bins y puede introducir estructura falsa en los
bordes. Con $40$--$80$ bins suele bastar; suba solo si necesita separar detalles
finos.
\end{truco}

\subsection{La regularización: $\log\alpha$ y atajos}

Solo para Histograma (y 2D):

\begin{itemize}[nosep]
  \item \menu{log$_{10}\alpha$}: fuerza de la regularización
        (§\ref{sec:regularizacion}). $\alpha$ es adimensional: $\log\alpha
        \approx 0$ es el equilibrio natural. Menos $\alpha$ $\to$ más detalle y
        más ruido; más $\alpha$ $\to$ más suave.
  \item Botones \boton{Fina} / \boton{Media} / \boton{Suave}: atajos que fijan
        $\log\alpha$ a valores progresivamente mayores (de más detalle a más
        suavidad), para tantear rápido.
  \item \boton{L-curve $\alpha$}: abre el diálogo que \textbf{elige $\alpha$ con
        criterio} (§\ref{sec:lcurve}). Es la vía recomendada.
\end{itemize}

\subsection{Controles de la 2D}

Aparecen solo con forma \textbf{2D}: \menu{mín/máx/bins del segundo eje} (p.\,ej.\
$\dEQ$) y \menu{log$_{10}\alpha$ del segundo eje}, análogos a los de la primera
variable pero para el eje $Y$. El botón \boton{Ver mapa 2D} abre el mapa
(§\ref{sec:dist-2d}).

\subsection{Otros}

\begin{itemize}[nosep]
  \item \menu{Añadir componentes nítidas activas}: ajusta a la vez la distribución
        y los componentes discretos activos (§\ref{sec:nitidos}).
  \item \menu{Cargar P fija…}: solo con forma \textbf{Fija}; lee una $P$ de
        fichero (dos columnas: valor y peso).
\end{itemize}

\section{Formas de la distribución}
\label{sec:formas-dist}

Fitbauer ofrece varias \textbf{formas} (\menu{Forma:}) para el $P(\BHF)$. Unas
son \emph{no paramétricas} (dejan libre el peso de cada bin) y otras
\emph{paramétricas} (imponen una función con pocos parámetros):

\begin{center}
\begin{tabular}{@{}lp{0.62\textwidth}@{}}
\toprule
\textbf{Forma} & \textbf{Descripción} \\
\midrule
\textbf{Histograma} & Inversión de Hesse--Rübartsch: pesos libres por bin con
regularización. La más general y flexible; necesita elegir $\alpha$. \\[3pt]
\textbf{Gaussiana} & Una única gaussiana en $\BHF$ (media y anchura). Rígida pero
muy estable. \\[3pt]
\textbf{VBF} & Suma de varias gaussianas (Voigt-Based Fitting, Rancourt--Ping);
guarda $A_k, \mu_k, \sigma_k$ por componente (§\ref{sec:vbf-areas}). \\[3pt]
\textbf{Binomial} & Distribución binomial (modelos de entornos de vecinos, p.\,ej.\
aleaciones). \\[3pt]
\textbf{Fija} & Distribución impuesta por el usuario (leída de fichero). \\[3pt]
\textbf{2D} & Distribución conjunta $P(\BHF, \dEQ)$ (§\ref{sec:dist-2d}). \\
\bottomrule
\end{tabular}
\end{center}

\begin{truco}
Elija \textbf{Gaussiana} o \textbf{VBF} cuando espere uno o pocos entornos bien
definidos (dan parámetros con significado directo); reserve el \textbf{Histograma}
para formas complejas o desconocidas, donde no quiera imponer una función.
\end{truco}

A continuación se detalla cada forma: qué modelo ajusta, qué controles usa y
cuándo elegirla.

\subsection{Histograma (Hesse--Rübartsch)}
\label{sec:forma-histograma}

Es la forma \textbf{no paramétrica} y la más general. No impone ninguna función:
trata el peso de \emph{cada bin} como una incógnita independiente. El espectro se
escribe como combinación lineal de los subespectros de la malla,
\[
  T(v) \;=\; b + s\,v \;-\; \sum_{j=1}^{N_{\text{bins}}} P_j\, S_j(v),
  \qquad P_j \ge 0,
\]
y ajustar significa \textbf{resolver los pesos} $\{P_j\}$ que mejor reproducen los
datos. Es un problema lineal con no-negatividad, pero \textbf{mal condicionado}
(muchos $\{P_j\}$ dan casi el mismo espectro), de ahí la necesidad de
\textbf{regularizar} (§\ref{sec:regularizacion}). El método se remonta a Hesse y
Rübartsch (1974).

\begin{itemize}[nosep]
  \item \textbf{Usa}: malla (mín/máx/bins), núcleo ($\delta$, $\dEQ$/BHF fijo,
        $\Gamma$), \textbf{regularizador} y $\log\alpha$, correlaciones
        $\kappa_\delta/\kappa_q$.
  \item \textbf{Cuándo}: cuando la forma de $P$ es \emph{desconocida} o compleja
        (varios entornos, colas asimétricas) y no quiere imponer una función.
  \item \textbf{Precio}: hay que elegir bien $\alpha$ (use la L-curve,
        §\ref{sec:lcurve}); es la única forma que la necesita.
\end{itemize}

\subsection{Gaussiana}
\label{sec:forma-gaussiana}

Forma \textbf{paramétrica} mínima: $P$ es \emph{una única gaussiana} sobre la
malla. Ajusta su \textbf{media} $\mu$ (el campo central) y su \textbf{anchura}
$\sigma$ (más la amplitud y la línea base). Solo dos parámetros de forma, así que
es \textbf{muy estable} y no sufre mal condicionamiento: no usa regularizador ni
$\alpha$. Internamente es exactamente el \textbf{VBF con una sola componente}
($N=1$, §\ref{sec:forma-vbf} y anexo~\ref{ap:vbf}): comparten algoritmo, y elegir
«Gaussiana» equivale a fijar $N=1$ con línea Lorentziana.

\begin{itemize}[nosep]
  \item \textbf{Cuándo}: hay \emph{un} entorno ensanchado (una fase con dispersión
        de campos aproximadamente simétrica), o como \textbf{primera
        aproximación} rápida antes de complicar el modelo.
  \item \textbf{Da}: $\mu$ y $\sigma$ con significado físico directo.
\end{itemize}

\subsection{VBF (Voigt-Based Fitting, Rancourt--Ping)}
\label{sec:forma-vbf}

Generaliza la Gaussiana a una \textbf{suma de $N$ gaussianas}
(\menu{Componentes VBF}, $N=1\ldots6$). Cada componente $k$ tiene amplitud $A_k$,
media $\mu_k$ y anchura $\sigma_k$, todas ajustables:
\[
  P(\BHF) \;=\; \sum_{k=1}^{N} A_k \,
  \mathcal{N}(\BHF;\,\mu_k,\sigma_k).
\]
Es el método de Rancourt y Ping (1991). Su gran ventaja frente al Histograma:
entrega \textbf{parámetros con significado} por subpoblación (posición, anchura y
\textbf{área}, §\ref{sec:vbf-areas}) en vez de un perfil sin estructura, y es más
estable (pocos parámetros, sin regularización). El anexo~\ref{ap:vbf} desarrolla
su matemática (modelo directo, perfil Voigt, parametrización, poblaciones y el
caso $N=1$ que corresponde a la forma \menu{Gaussiana}).

\begin{itemize}[nosep]
  \item \textbf{Cuándo}: espera \textbf{2--3 subpoblaciones} magnéticas bien
        definidas y quiere cuantificar su \emph{proporción}.
  \item \textbf{Consejo}: empiece con $N=1$--$2$ y suba solo si el residuo lo pide;
        más gaussianas de las necesarias se solapan y quedan mal definidas.
  \item \textbf{Admite}: correlaciones $\kappa_\delta/\kappa_q$.
\end{itemize}

\subsection{Binomial}
\label{sec:forma-binomial}

Forma \textbf{paramétrica} con una sola incógnita de forma: $P$ sigue una
\textbf{distribución binomial} $B(n,p)$ sobre la malla (con $n = N_{\text{bins}}-1$),
y se ajusta la \textbf{probabilidad} $p\in(0,1)$ (la media cae en $p\,n$).

Su motivación es física: en una \textbf{aleación}, el campo hiperfino de un átomo
de Fe depende de \textbf{cuántos átomos de soluto} tiene en sus capas de vecinos.
Si el soluto se distribuye al azar, ese número sigue una binomial, y por tanto la
distribución de campos también. El parámetro $p$ se relaciona con la
\textbf{concentración} de soluto.

\begin{itemize}[nosep]
  \item \textbf{Cuándo}: sistemas de \emph{vecinos} (aleaciones diluidas) donde
        cada configuración de vecinos da un campo discreto.
\end{itemize}

\subsection{Fija}
\label{sec:forma-fija}

Usa una distribución $P$ \textbf{que usted proporciona} en un fichero de dos
columnas (valor y peso), cargado con \menu{Cargar P fija…}. La forma \emph{no} se
ajusta (solo amplitud y línea base): sirve para \textbf{reutilizar} una $P$
conocida —de un ajuste anterior, de la literatura o teórica— y comparar o simular
con ella sobre distintos espectros.

\subsection{2D: dos variables a la vez}
\label{sec:forma-2d}

Distribuye \textbf{dos} magnitudes conjuntamente (p.\,ej.\ $P(\BHF,\dEQ)$) sobre
una malla bidimensional. Es el caso más general y más costoso; se trata aparte en
§\ref{sec:dist-2d}.

\subsection{Áreas de cada gaussiana (VBF)}
\label{sec:vbf-areas}

En modo \textbf{VBF} con más de una gaussiana, cada componente aporta un área
$A_k$ (su integral en la malla, porque el núcleo gaussiano está normalizado a
área unidad). Fitbauer informa de ese \textbf{área} y de su \textbf{porcentaje
relativo} $A_k/\sum_j A_j$, que representa la \textbf{población de cada entorno}:

\begin{itemize}[nosep]
  \item En el \textbf{panel de estado}: claves \ruta{vbf$k$\_area},
        \ruta{vbf$k$\_area\%} y \ruta{vbf$k$\_mu} por componente.
  \item En el \textbf{resumen del ajuste} y en los \textbf{informes} (tabla con
        $A$, área \%, $\mu$ y $\sigma$).
\end{itemize}

Esto convierte al VBF en una vía cómoda para cuantificar la proporción de dos o
más subpoblaciones magnéticas sin recurrir a un histograma.

\section{Regularización del histograma}
\label{sec:regularizacion}

Ajustar un \textbf{Histograma} a datos reales es un problema \emph{mal
condicionado}: muchos $P(\BHF)$ distintos reproducen casi igual de bien el
espectro, y sin control aparecen oscilaciones espurias. La solución es
\textbf{regularizar}: añadir al ajuste un término que penaliza las soluciones
poco suaves, controlado por un parámetro $\alpha$.

\subsection{Regularizadores}

\begin{center}
\begin{tabular}{@{}lp{0.6\textwidth}@{}}
\toprule
\textbf{Regularizador} & \textbf{Efecto / cuándo usarlo} \\
\midrule
\textbf{Tikhonov} & Penaliza la curvatura (segunda diferencia): soluciones
\emph{suaves}. Opción por defecto, adecuada para distribuciones anchas y
redondeadas. \\[4pt]
\textbf{TV} (variación total) & Penaliza la variación total: preserva
\emph{bordes} y escalones; útil si sospecha poblaciones bien separadas. \\[4pt]
\textbf{MaxEnt} (máxima entropía) & Busca la distribución más «uniforme»
compatible con los datos (L-BFGS-B con gradiente analítico); evita introducir
estructura no justificada. \\
\bottomrule
\end{tabular}
\end{center}

\subsection{El parámetro $\alpha$}

$\alpha$ equilibra \textbf{fidelidad a los datos} (poco $\alpha$ $\to$ sigue el
ruido, distribución rugosa) frente a \textbf{suavidad} (mucho $\alpha$ $\to$
distribución lisa pero que ajusta peor). En Fitbauer $\alpha$ es
\textbf{adimensional}: $\log_{10}\alpha \approx 0$ es el equilibrio natural para
cualquier espectro, y los valores útiles suelen caer entre $-1$ y $+2$.

\begin{aviso}
$\alpha$ es la elección más delicada del ajuste por histograma. No lo deje al
azar: use la \textbf{L-curve} (§\ref{sec:lcurve}) para elegirlo con criterio.
\end{aviso}

\section{La L-curve: elegir $\alpha$ con criterio}
\label{sec:lcurve}

Elegir $\alpha$ «a ojo» es arriesgado: un valor demasiado bajo mete ruido en la
distribución y otro demasiado alto la aplana y esconde estructura real. La
\textbf{L-curve} da un criterio objetivo. Es la herramienta más importante del
ajuste por histograma.

\subsection{Qué es y por qué tiene forma de L}

Para cada $\alpha$ del barrido, el ajuste produce una distribución $P_\alpha$, y de
ella se miden \textbf{dos cantidades} enfrentadas:

\begin{itemize}[nosep]
  \item \textbf{Desajuste} (residuo) $\lVert K P_\alpha - y\rVert$: cuánto se
        aparta el modelo de los datos. Queremos que sea \emph{pequeño}.
  \item \textbf{Rugosidad} $\lVert L\,P_\alpha\rVert$: cuánto «tiembla» la
        distribución (la norma del operador de suavidad $L$). También queremos
        que sea \emph{pequeña}.
\end{itemize}

No se pueden minimizar las dos a la vez: bajar $\alpha$ mejora el ajuste pero
ensucia la distribución, y subirlo hace lo contrario. Si se dibuja
\textbf{rugosidad frente a desajuste} en escala log-log al variar $\alpha$, sale
una curva con forma de \textbf{L}:

\begin{itemize}[nosep]
  \item \textbf{Rama vertical} ($\alpha$ pequeño, \emph{sub-regularizado}):
        seguir bajando $\alpha$ apenas mejora el ajuste, pero la rugosidad se
        \textbf{dispara} —la solución empieza a ajustar el ruido.
  \item \textbf{Rama horizontal} ($\alpha$ grande, \emph{sobre-regularizado}):
        la distribución ya es lisa; subir más $\alpha$ apenas la suaviza, pero el
        ajuste \textbf{empeora} —se pierde estructura real.
  \item \textbf{La esquina}: el punto donde la curva gira. Es el
        \textbf{mejor compromiso}: la $\alpha$ más pequeña posible que todavía
        mantiene la distribución razonablemente suave.
\end{itemize}

Formalmente, la esquina es el punto de \textbf{máxima curvatura} de la L
(criterio de Hansen). Fitbauer la detecta automáticamente (curvatura de Menger
sobre tríos de puntos) y la marca en \textcolor{fbgreen}{verde} como
\emph{sugerencia}.

\subsection{El diálogo, panel a panel}

\begin{itemize}[nosep]
  \item \textbf{Izquierda — la L}: rugosidad frente a desajuste (log-log), con
        cada punto etiquetado por su $\log\alpha$. La esquina automática en verde
        y el punto \textbf{elegido} en \textcolor{fbred}{naranja}.
  \item \textbf{Derecha — $\alpha$ vs RMS}: el error de ajuste (RMS) frente a
        $\alpha$; ayuda a ver dónde el ajuste empieza a degradarse.
\end{itemize}

Ambas figuras están \textbf{acopladas}: al hacer zoom en una, la otra se limita al
mismo rango de $\alpha$. La barra de herramientas permite ampliar la zona de la
esquina.

\guicap{gui_lcurve}{Diálogo de la L-curve: a la izquierda la curva
rugosidad$\leftrightarrow$desajuste con la esquina automática (verde) y el punto
elegido (naranja); a la derecha $\alpha$ frente a RMS. Un clic elige $\alpha$.}

\subsection{Elegir $\alpha$: automático o a mano}

La detección automática de la esquina \textbf{no siempre acierta} (curvas ruidosas,
esquinas poco marcadas), así que el diálogo es \textbf{interactivo}:

\begin{itemize}[nosep]
  \item \textbf{Haga clic} sobre cualquiera de las dos figuras para \textbf{elegir}
        el $\alpha$ escaneado más cercano al clic; el punto se resalta en ambas.
  \item El botón \boton{Usar $\alpha$=…} aplica el valor \emph{elegido} (arranca en
        la sugerencia automática) al ajuste y cierra el diálogo.
\end{itemize}

\subsection{Cómo interpretarla en la práctica}

\begin{enumerate}[nosep]
  \item Acepte la \textbf{esquina automática} como primera aproximación y observe
        la $P(\BHF)$ resultante.
  \item Si la ve \textbf{demasiado rugosa} (picos que parecen ruido), suba un poco
        $\alpha$ (pinche un punto a la derecha en la L). Si la ve
        \textbf{demasiado lavada} (ha perdido detalle que el residuo delataba),
        baje $\alpha$ (un punto a la izquierda).
  \item Compruebe la \textbf{estabilidad}: la solución correcta cambia \emph{poco}
        al mover $\alpha$ un paso arriba o abajo. Si cambia mucho, está en una
        rama, no en la esquina.
\end{enumerate}

\begin{aviso}
A veces \textbf{no hay una esquina clara} (la L es redondeada). Suele significar
que los datos no exigen mucha estructura: elija una $\alpha$ de la zona de
transición y valídela por la estabilidad de $P(\BHF)$ y por el $\chi^2_r$ del
ajuste, no por la esquina. Y recuerde: la L-curve solo aplica al \textbf{Histograma}
(las formas paramétricas no usan $\alpha$).
\end{aviso}

\section{Correlación $\delta(H)$ y $\dEQ(H)$}
\label{sec:correlacion}

Hasta aquí hemos supuesto que $\delta$ y $\dEQ$ son \textbf{iguales para todos}
los subespectros de la distribución, y solo variaba el campo. Pero físicamente eso
no siempre es cierto: el desorden que dispersa el campo suele dispersar
\emph{también} $\delta$ y $\dEQ$, y de forma \textbf{correlacionada}. Por ejemplo,
un átomo de Fe con más vecinos de un cierto tipo tiene a la vez un campo distinto
\emph{y} una densidad electrónica distinta (otro $\delta$): las dos magnitudes se
mueven juntas.

Fitbauer captura esa correlación con dos \textbf{pendientes} (\emph{opt-in}, cero =
comportamiento clásico sin correlación):

\begin{itemize}[nosep]
  \item $\kappa_\delta = d\delta/dH$: hace que $\delta(H) = \delta_0 +
        \kappa_\delta\,(H-H_{\text{ref}})$, es decir, el desplazamiento isomérico
        \emph{varía linealmente} con el campo del subespectro.
  \item $\kappa_q = d\dEQ/dH$: lo mismo para el cuadrupolo.
\end{itemize}

\textbf{Qué implica.} Con $\kappa_\delta \neq 0$, los subespectros de campo alto
usan un $\delta$ distinto que los de campo bajo, de modo que la distribución ya no
solo «abre» las líneas sino que las \emph{desplaza} de forma dependiente del
campo. Es la manera de modelar, con \emph{un} solo parámetro extra, que el entorno
químico y el magnético cambian a la par. Si su material lo requiere, verá que el
residuo mejora al liberar $\kappa_\delta$/$\kappa_q$ (desmarcando su candado).

Aplican a las formas \textbf{Histograma} y \textbf{VBF}. Los valores ajustados de
$\kappa_\delta$ y $\kappa_q$ se muestran en el panel de estado y en los informes.

\begin{truco}
Empiece siempre con $\kappa=0$. Actívelos solo si, tras un buen ajuste sin
correlación, el residuo muestra un patrón sistemático que sugiere que $\delta$ o
$\dEQ$ dependen del campo. Un $\kappa$ mal usado puede compensar artificialmente
otros defectos del modelo.
\end{truco}

\section{Distribuciones 2D: $P(\BHF, \dEQ)$}
\label{sec:dist-2d}

Cuando campo y cuadrupolo varían de forma \textbf{no correlacionable con una
recta} (es decir, cuando $\kappa_\delta/\kappa_q$ no bastan), la forma \textbf{2D}
ajusta una densidad conjunta $P(\BHF, \dEQ)$ sobre una malla bidimensional. Cada
eje tiene sus propios controles:

\begin{itemize}[nosep]
  \item Eje $X$ (p.\,ej.\ $\BHF$): \menu{mín/máx/bins} y \menu{log$_{10}\alpha$}
        habituales.
  \item Eje $Y$ (p.\,ej.\ $\dEQ$): \menu{mín/máx/bins del segundo eje} y su propio
        \menu{log$_{10}\alpha$}, que se \textbf{regulariza por separado}.
\end{itemize}

El resultado se muestra como un \textbf{mapa} (\emph{heatmap} con contornos)
—accesible con \boton{Ver mapa 2D}— además de las distribuciones marginales de
cada variable. El desarrollo matemático del caso bidimensional está en el
anexo~\ref{ap:dist-2d}.

\begin{aviso}
El 2D tiene \textbf{muchos más incógnitas} (una por celda de la malla) que el 1D,
así que es más costoso y más difícil de regularizar. Empiece con \textbf{pocos
bins} en cada eje (p.\,ej.\ $15\times15$) para tantear, compruebe que el mapa es
estable, y solo entonces afine la malla. Dos $\alpha$ demasiado bajos producen un
mapa lleno de picos falsos.
\end{aviso}

\section{Componentes nítidos junto a la distribución}
\label{sec:nitidos}

Es habitual que una muestra combine una \textbf{fase cristalina} (un sexteto
nítido) con una \textbf{fase distribuida}. Con \menu{P(BHF): sumar componentes
activos nítidos} se ajustan \textbf{simultáneamente} uno o varios componentes
discretos y la distribución. El gráfico dibuja entonces cada nítido por separado,
la \textbf{envolvente} de la distribución y el modelo total, de modo que se ve la
contribución de cada fase.

\section{Paso a paso: ajustar una \texorpdfstring{$P(\BHF)$}{P(BHF)} desde cero}
\label{sec:paso-a-paso}

Si es la primera vez, siga esta receta. Cada paso remite a la sección que lo
explica en detalle.

\begin{enumerate}[nosep]
  \item \textbf{Cargue y calibre.} Abra el espectro y asegúrese de tener una
        calibración fijada (capítulo~\ref{ch:calibracion}); sin la escala de
        velocidad correcta, los campos saldrán sesgados.
  \item \textbf{Entre en modo distribución} y elija la \textbf{variable}
        $P(\BHF)$ (§\ref{sec:dist-variable}).
  \item \textbf{Tantee con una forma sencilla.} Empiece con \menu{Gaussiana} (o
        \menu{VBF} con 1--2 componentes): son estables y le dan una idea rápida de
        dónde está el campo y cuánto de ancho es. Fije $\delta$ al de su fase y una
        $\Gamma \approx 0{,}30$~mm/s (§\ref{sec:panel-kernel}).
  \item \textbf{Ajuste el rango de la malla} (mín/máx) para que cubra donde hay
        campo, sin sobrar (§\ref{sec:panel-dist}); use $\sim 50$ bins.
  \item \textbf{Lance el ajuste} (\menu{Ajuste > Ajustar}). Mire la $P(\BHF)$ y el
        residuo: ¿reproduce las líneas?
  \item \textbf{¿Forma compleja?} Si una o pocas gaussianas no bastan, pase a
        \menu{Histograma} (§\ref{sec:formas-dist}), que no impone una función.
  \item \textbf{Elija $\alpha$ con la L-curve} (§\ref{sec:lcurve}): no lo deje por
        defecto. Acepte la esquina o pinche un punto y compare la $P(\BHF)$.
  \item \textbf{Afine el modelo} si procede: active \textbf{componentes nítidos}
        para una fase cristalina superpuesta (§\ref{sec:nitidos}), o las
        pendientes $\kappa_\delta/\kappa_q$ si $\delta$/$\dEQ$ dependen del campo
        (§\ref{sec:correlacion}).
  \item \textbf{Documente.} Con VBF, anote las \textbf{áreas} de cada gaussiana
        (proporción de fases, §\ref{sec:vbf-areas}). Guarde la \textbf{sesión} y
        exporte el \textbf{informe} con la tabla de la distribución.
\end{enumerate}

\begin{truco}
Regla de oro: vaya \textbf{de lo simple a lo complejo}. Una Gaussiana estable que
explica el 95\% del espectro es mejor punto de partida que un Histograma mal
regularizado. Suba en complejidad (VBF $\to$ Histograma $\to$ 2D) solo cuando el
residuo lo justifique, y vigile siempre que la solución sea \textbf{estable} ante
pequeños cambios de $\alpha$, $\Gamma$ o número de bins.
\end{truco}
